%!TEX root = ./main.tex

\section[Wrap-up]{The Design Space}

% SECTION TIMING: 72:00--75:00 (3 frames). Pure synthesis: no new papers, no new
% numbers, no new claims. Everything here echoes section 01 and the three acts.
% If Act III ran long, drop frame 2 ("How to read a paper like today's") and speak
% its three lines over the final frame instead.
% LAYOUT RULE: every bullet must fit on one line. Shorten the text, do not let it wrap.

% TEACHING NOTE (1:15): read each row as a chant -- point left, say the assumption,
% point right, say the move. The left column echoes, keyword for keyword, the five
% assumptions from section 01. Students heard all five already; this frame only lines
% them up so the pattern is visible at once. Do not re-explain any paper here.
% Close by saying: the diagram did not change today -- what changed is how much of
% it you take on faith.
\begin{frame}{Five assumptions walked in. None walked out.}
  \vspace{0.3em}
  \centering
  \small
  \renewcommand{\arraystretch}{1.45}
  \begin{tabular}{@{}l@{\hspace{0.6em}}l@{\hspace{0.6em}}l@{}}
    \toprule
    & \textbf{The move} & \textbf{Act} \\
    \midrule
    \textbf{On the ground}
      & Earth \(\to\) orbit
      & \textcolor{ranblue}{\textbf{I}} \\
    \textbf{Only a radio}
      & GPU, AI, antivirus moved in
      & \textcolor{ranpurple}{\textbf{II}}--\textcolor{attackred}{\textbf{III}} \\
    \textbf{Transmit, then compute}
      & the wave computes in flight
      & \textcolor{ranpurple}{\textbf{II}} \\
    \textbf{Inside links trusted}
      & fronthaul is an attack surface
      & \textcolor{attackred}{\textbf{III}} \\
    \textbf{Fake base station}
      & perfect timing instead
      & \textcolor{attackred}{\textbf{III}} \\
    \bottomrule
  \end{tabular}

  \vspace{0.35em}
  {\scriptsize\textcolor{softgray}{...and one assumption was never even a box:
   \emph{the cache sits on the ground} (Act I).}\par}

  \vspace{0.35em}
  \normalsize
  One synthesis: \textbf{openness} \(\to\) innovation \textbf{and} attack surface.

  \vspace{0.35em}
  \takeaway{Seven papers and one teaser, one shared trick:\\
    find the sentence the diagram never says.}
\end{frame}

% TEACHING NOTE (0:45): deliver as three quick beats, fifteen seconds each. This is
% the recipe students will reuse in the practicum, so say it like a recipe, not a sermon.
\begin{frame}{How to read a paper like today's}
  \vspace{0.8em}
  \large
  \begin{enumerate}\itemsep14pt
    \item Find the \textbf{assumption} it breaks.\\
      {\small\textcolor{softgray}{Every paper today had exactly one.}}
    \item Ask what it \textbf{deliberately does not solve}.\\
      {\small\textcolor{softgray}{That silence is where the next paper lives.}}
    \item Ask what \textbf{you} would build on top of it.\\
      {\small\textcolor{softgray}{That question is the whole practicum.}}
  \end{enumerate}

  \vspace{0.7em}
  \normalsize
  Steps 1 and 2 need the paper. Step 3 needs \textbf{you}.
\end{frame}

% TEACHING NOTE (1:00): let the big line sit in silence for a few seconds before
% speaking. Then the gear-shift: Classes 2--4 taught how it works; from here the
% course asks what the students would redesign. Name Sni5Gect last, as the bridge
% to the practicum, and end there.
\begin{frame}[plain]
  \vspace{2.2em}
  \centering
  {\huge The 5G architecture diagram\\[0.35em] is \textbf{not the answer}.\par}

  \vspace{1.1em}
  {\huge It is \textcolor{coreorange}{\textbf{the design space}}.\par}

  \vspace{2.2em}
  \normalsize\textcolor{softgray}{%
    Until today, this course taught you how it works.\\[0.25em]
    From here, it asks what \textit{you} would redesign.\\[0.25em]
    The paper practicum starts soon --- and Sni5Gect is already on the list.}
\end{frame}
